\documentclass[11pt]{article}
\usepackage[margin=1in]{geometry}
\usepackage[T1]{fontenc}
\usepackage{lmodern,microtype,enumitem,xurl,hyperref}
\hypersetup{colorlinks=true,urlcolor=blue,linkcolor=blue}
\setlength{\parindent}{0pt}
\setlength{\parskip}{5pt}
\setlength{\emergencystretch}{3em}
\title{S-two recursive privacy blinding: disclosure summary}
\author{}
\date{September 15, 2026}
\begin{document}
\maketitle

\textbf{Status:} Technical summary of the saved audit report, prepared for
disclosure review. The original report remains in the research archive.
This version summarizes evidence, impact, limitations, and remediation;
it does not include the witness-recovery procedure or executable artifacts.

\textbf{Audited source:}
\nolinkurl{starkware-libs/proving@cd7bc5f4697fb188a27e09f9242f1dd76df8afdc}.

\textbf{Affected area:} The recursive privacy prover and its shared
\texttt{add\_zk\_blinding} helper.

\section*{Finding}

The saved audit reports a source-compatible valid toy circuit for which
the disclosed blinding transcript determines a private witness on a
specified class of public opening configurations. The circuit admits
different private witnesses with the same public output. The reported
algebraic result holds even when the masks are fresh and independent.

The issue concerns whether the helper can be used as a generic
witness-hiding transformation. Adding random rows does not alone prove
that the complete collection of disclosed evaluations hides the private
part of a trace. The relevant condition concerns their combined
information content and the way the masking randomness enters them.

\section*{Evidence and what it establishes}

The original report records exact finite-field calculations showing
that the public observation map distinguishes valid private witnesses
in its stated toy circuit. It also records checks of finalization and
padding compatibility. These are reported audit results; preparation
of this summary did not rerun the recovery experiment or a native
recursive verifier.

The audit separately analyzes opening configurations under an idealized
challenge model. That calculation is not a theorem about the complete
Fiat--Shamir transcript distribution of the deployed application.
An exact conditional algebraic finding and a concrete end-to-end
privacy claim require different evidence.

\section*{Impact and unresolved scope}

The finding challenges an unrestricted claim that this helper hides
witnesses for arbitrary valid circuits under the audited transcript
interface. A privacy argument for the fixed recursive application must
therefore account for its actual witness relation, masking distribution,
and complete transcript.

The saved report does not establish witness recovery from an accepted
production privacy transaction. Two material questions remain:
whether the demonstrated relation is admitted by the fixed recursive
verifier, and whether the relevant opening event has a proved probability
under the complete Fiat--Shamir process. No deployed secret or live
service was accessed in the reported audit.

The report also identifies a separate randomness obligation: the audited
path derives its masking seed from a proof commitment. Any privacy proof
must justify the resulting distribution and its dependence on the
witness; the mere presence of a seed is not such a justification.

\section*{Recommended remediation and validation}

\begin{enumerate}[leftmargin=*]
\item Specify the privacy property, supported witness relation, and
complete transcript, including query and out-of-domain disclosures.
\item Establish a simulator argument or an equivalent information-hiding
proof for that relation and transcript. Account for all disclosed
projections together when choosing masking dimensions and placement.
\item Give the prover-randomness interface an explicit entropy and
independence contract, and justify any deterministic seed derivation.
\item Add regression checks comparing valid witnesses with the same
public statement, together with transcript-level leakage checks and
proof/verifier round trips for supported component configurations.
\item Record affected versions, the implemented correction, and the
security argument for the corrected configuration.
\end{enumerate}

\section*{Relationship to the research paper}

This privacy finding is separate from the repository's list-decoding
and proximity-gap results. Those results concern quantitative coding
premises in soundness analyses. They do not establish the privacy
finding, and this finding does not establish an end-to-end soundness
failure or an improved coding-challenge submission.

\end{document}
